\subsection*{Justificación Fila 4}

El join por la derecha conserva todas las tuplas de la relación de la derecha ($S$), además de las coincidencias del join natural. Las tuplas de $S$ que no tienen coincidencia con $R$ se conservan y se rellenan con nulos (NULL) en los atributos correspondientes a $R$.

Las coincidencias naturales ya las calculamos en el paso anterior y sabemos que nos dan 5 tuplas. Revisamos las tuplas de $S$ que no hicieron match con $R$:
\begin{itemize}
    \item $(w, 3, 0)$ en $S$ (no hay valor $w$ en la columna $B$ de $R$).
\end{itemize}
Esta única tupla se conserva y se rellena con NULL para la columna $A$ de $R$.

Por lo tanto, la cardinalidad del resultado es $5 + 1 = \mathbf{6}$ \textbf{tuplas}.

\begin{center}
$(R \bowtie= S) =$ \quad
\begin{tabular}{|c|c|c|c|}
\hline
\textbf{A} & \textbf{B} & \textbf{C} & \textbf{D} \\ \hline
1 & x & 0 & 3 \\
3 & x & 0 & 3 \\
2 & y & 1 & 2 \\
2 & y & 3 & 3 \\
2 & y & 4 & 2 \\
\text{NULL} & w & 3 & 0 \\ \hline
\end{tabular}
\end{center}
